\documentclass[11pt, tikz, border=12pt]{standalone}
\usepackage[fontset=fandol]{ctex}
\usetikzlibrary{positioning, shadows, arrows.meta, backgrounds, calc, fit}

\definecolor{cBase}{RGB}{120, 138, 148}
\definecolor{cEvo}{RGB}{136, 84, 162}
\definecolor{cOpt}{RGB}{204, 112, 34}
\definecolor{cHum}{RGB}{52, 126, 186}
\definecolor{cOut}{RGB}{46, 142, 84}

\newcommand{\cardcontent}[5]{%
    \begin{minipage}[c]{1.05cm}
        \centering
        \tikz[baseline=-0.5ex]{
            \fill[#1, rounded corners=3pt] (-0.42,-0.42) rectangle (0.42,0.42);
            \begin{scope}[draw=white, thick, line cap=round, line join=round, scale=1.25]
                #3
            \end{scope}
        }
    \end{minipage}%
    \hspace{1.8mm}%
    \begin{minipage}[c]{#2}
        \vspace{2pt}
        \textbf{\fontsize{9.8pt}{12pt}\selectfont\textcolor{black!88}{#4}}\\[3pt]
        \parbox{#2}{\fontsize{8.2pt}{11.5pt}\selectfont\color{black!68}#5}
        \vspace{2pt}
    \end{minipage}%
}

\begin{document}
\begin{tikzpicture}[
    font=\sffamily,
    card/.style={draw=black!18, fill=white, thick, rounded corners=5pt, inner sep=8pt},
    arrow/.style={-{Latex[scale=1.12]}, thick, draw=black!50},
    optarrow/.style={-{Latex[scale=1.08]}, thick, draw=cOpt!90!black},
    humarrow/.style={-{Latex[scale=1.08]}, thick, draw=cHum!90!black},
    edgelabel/.style={font=\scriptsize\bfseries, fill=white, inner sep=2.5pt, rounded corners=2pt}
]

    % ---- 顶部：基础输入 ----
    \node[card] (input) at (0, 0) {
        \cardcontent{cBase}{6.2cm}{
            \draw (-0.12,-0.15) rectangle (0.12,0.15);
            \draw (-0.05, 0.05) -- (0.05, 0.05);
            \draw (-0.05,-0.05) -- (0.05,-0.05);
        }{基础指令--回答对 $(x, y)$（来自图遍历生成）}{包含单一约束条件与基础推理链条。}
    };

    % ---- 中央三步演化 ----
    \node[card] (step1) at (0, -3.9) {
        \cardcontent{cEvo}{6.2cm}{
            \draw (0,-0.05) circle(0.06);
            \draw[->] (0.05,-0.01) -- (0.15, 0.12);
            \draw[->] (0.06,-0.05) -- (0.18,-0.05);
            \draw[->] (0.05,-0.09) -- (0.15,-0.22);
        }{第一步：操作选择与策略生成}{分析任务类型、推理步数与所属数据类型，从五类演化操作集合 $\mathcal{O}$ 中选择适配操作及强度参数 $\theta$。}
    };

    \node[card] (step2) at (0, -7.6) {
        \cardcontent{cEvo}{6.2cm}{
            \draw[->] (-0.12,0.05) arc (180:20:0.12);
            \draw[<-] (-0.08,-0.05) arc (180:340:0.12);
        }{第二步：操作执行与答案同步重写}{依次执行演化操作生成 $x'$，并以 $x'$ 为输入重新生成完整回答 $y'$，确保推理深度与新增约束严格对应。}
    };

    \node[card] (step3) at (0, -11.3) {
        \cardcontent{cEvo}{6.2cm}{
            \draw (0,-0.18) -- (0.15,0) -- (0.15,0.18) -- (-0.15,0.18) -- (-0.15,0) -- cycle;
            \draw (-0.05,-0.02) -- (-0.01,-0.06) -- (0.08,0.03);
        }{第三步：三项内嵌自检与过滤}{（1）复杂度提升检验；（2）约束组合可解性检验（排查互斥条件）；（3）回答完整性检验。}
    };

    % ---- 底部：成功输出 ----
    \node[card] (output) at (0, -15.3) {
        \cardcontent{cOut}{6.2cm}{
            \draw (-0.15,0.08) arc (180:360:0.15 and 0.05);
            \draw (-0.15,0.08) arc (180:0:0.15 and 0.05);
            \draw (-0.15,-0.04) arc (180:360:0.15 and 0.05);
            \draw (-0.15,-0.16) arc (180:360:0.15 and 0.05);
            \draw (-0.15,0.08) -- (-0.15,-0.16);
            \draw (0.15,0.08) -- (0.15,-0.16);
        }{复杂候选 SFT 指令变体 $(x', y')$}{进入 3.6 节自动化多维质量筛选模块。}
    };

    % ---- 右侧反馈环：Optimizer LLM ----
    \node[card] (optimizer) at (8.4, -7.6) {
        \cardcontent{cOpt}{4.1cm}{
            \draw (-0.1,-0.1) -- (0.05,0.05);
            \draw (0.05,0.05) arc (225:45:0.08);
            \fill (0.16, 0.16) circle(0.03);
        }{Optimizer LLM}{统计批次失败模式，迭代修订演化操作提示词约束。}
    };

    % ---- 左侧反馈环：人工轻量介入 ----
    \node[card] (human) at (-8.4, -7.6) {
        \cardcontent{cHum}{4.1cm}{
            \fill (0,0.1) circle (0.07);
            \draw (-0.18,-0.15) .. controls (-0.18,-0.02) and (0.18,-0.02) .. (0.18,-0.15);
        }{人工轻量级抽检}{启动期审阅 100 条 + 常规批次 5\% 随机抽检监控。}
    };

    % ---- Evolution LLM 虚线包络框（自动适配 step1..step3，绝不压框） ----
    \begin{scope}[on background layer]
        \coordinate (evoTop) at ([yshift=0.52cm]step1.north);
        \coordinate (evoBot) at ([yshift=-0.18cm]step3.south);
        \node[fit=(step1)(step2)(step3)(evoTop)(evoBot), draw=cEvo!45, fill=cEvo!4, dashed, thick, rounded corners=6pt, inner xsep=8pt, inner ysep=4pt] (evobox) {};
        \node[font=\footnotesize\bfseries, text=cEvo!85!black, anchor=north west] at ([xshift=8pt, yshift=-4pt]evobox.north west) {Evolution LLM 自主演化主干};
    \end{scope}

    % ---- 中央纵向主干箭头 ----
    \draw[arrow] (input.south) -- (step1.north);
    \draw[arrow] (step1.south) -- (step2.north);
    \draw[arrow] (step2.south) -- (step3.north);
    \draw[arrow] (step3.south) -- node[midway, right=3pt, font=\scriptsize\bfseries, text=cOut!90!black] {通过自检} (output.north);

    % ---- 右侧闭环反馈线（标签置于虚线框外侧走廊中央，避免虚线切字） ----
    \coordinate (optMidX) at ($(evobox.east)!0.5!(optimizer.west)$);
    \draw[optarrow, rounded corners=5pt]
        (step3.east) -- (optimizer.center |- step3.east)
        node[pos=0.58, edgelabel, text=cOpt!90!black] {未通过自检样本}
        -- (optimizer.south);

    \draw[optarrow, rounded corners=5pt]
        (optimizer.north) -- (optimizer.center |- step1.east)
        -- node[pos=0.42, edgelabel, text=cOpt!90!black] {更新提示词约束} (step1.east);

    % ---- 左侧人工监控线 ----
    \draw[humarrow, rounded corners=5pt]
        (output.west) -- (human.center |- output.west)
        node[pos=0.55, edgelabel, text=cHum!90!black] {5\% 随机抽检}
        -- (human.south);

    \draw[humarrow, rounded corners=5pt]
        (human.north) -- (human.center |- step1.west)
        -- node[pos=0.42, edgelabel, text=cHum!90!black] {宏观演化方向校准} (step1.west);

\end{tikzpicture}
\end{document}
